\begin{frame}{Ergebnisvideo und Blur-Qualität}
\begin{itemize}
    \item Ovaler Gaussian Blur wirkt natürlich und deckt Gesichter weich ab.
    \item Pixelate ist sichtbarer und kann bei kleinen Gesichtern härter wirken.
    \item Die Blur-Qualität hängt direkt am Detektions-Recall.
    \item Smoke-Videos zeigen die Pipeline, ersetzen aber keine finale Datenschutzprüfung.
\end{itemize}
\vspace{0.6em}
\begin{mynote}[Qualitätscheck]
Technisch gute Metriken reichen nicht: Am Ende muss das Video visuell auf verpasste Gesichter geprüft werden.
\end{mynote}
\end{frame}
